\documentclass[12pt]{article}
\usepackage[a4paper, margin=1in]{geometry}
\usepackage{amsmath,amssymb,amsthm}
\usepackage{mathtools}
\usepackage{mathrsfs}
\usepackage{enumitem}
\usepackage{xcolor}
\usepackage[hidelinks]{hyperref}
\usepackage{booktabs}
\usepackage{array}
\usepackage{fontspec}

\defaultfontfeatures{Ligatures=TeX}
\IfFontExistsTF{Times New Roman}
  {\setmainfont{Times New Roman}}
  {\setmainfont{TeX Gyre Termes}}
\IfFontExistsTF{Menlo}
  {\setmonofont{Menlo}}
  {\setmonofont{Latin Modern Mono}}

\newcolumntype{P}[1]{>{\raggedright\arraybackslash}p{#1}}

\newcommand{\tightlist}{%
  \setlength{\itemsep}{0pt}\setlength{\parskip}{0pt}}

\title{%
  \textbf{The Imprint Chain}\\
  \textbf{A Roadmap and Epistemic Grading of a Research Programme}%
}

\author{
  Sidong Liu, PhD \\
  \small iBioStratix Ltd \\
  \small \texttt{sidongliu@hotmail.com}
}

\date{May 2026\\
\small Zenodo DOI: \href{https://doi.org/10.5281/zenodo.20151594}{10.5281/zenodo.20151594}}

\begin{document}
\emergencystretch=2em
\raggedbottom
\hbadness=5000
\vbadness=5000

\maketitle

\begin{abstract}
This paper presents a roadmap rather than a new theorem. Its purpose is to display, grade, and mark the seams of a chain that runs from algebraic emergence through scale-stable imprinting and accessible world-formation to self-reading mind and the structural horizon of self-description. The central claim is that the chain is real as a research architecture, but its segments do not all have the same epistemic status.

The paper's two practical contributions are an explicit epistemic grading map and a set of reading paths for external readers who do not want to traverse the full body of work in chronological order.

\medskip

\noindent\textbf{Keywords}: emergence, imprint, self-reading, roadmap, programme-level bridge, epistemic grading, consciousness
\end{abstract}

\clearpage
\tableofcontents
\clearpage

\section{Introduction: Why a Roadmap Is Needed}

The body of work linking the Genesis Trilogy, the Consciousness Series, and \emph{The World as Imprint} now contains enough internal structure that external readers face an avoidable difficulty. The difficulty is not primarily one of disagreement. It is one of entry. Which claims are supposed to be read as formal results? Which are only programme-level bridges? Which are interface proposals between technical and philosophical languages? Which are structural resonances that orient interpretation but do not function as evidence?

This paper is written to answer those questions directly. It does not add a new mechanism to the chain from emergence to self-reading mind. It displays the chain, grades its segments, and marks where different kinds of claim meet.

The need for such a map became especially clear once \emph{The World as Imprint} introduced a compressed philosophical vocabulary capable of speaking across multiple registers at once \cite{WorldImprint}. That paper successfully distinguished universal imprint-carrying from uneven self-reading, but precisely because it was portable and programmatic, it left open a structural question for readers approaching from outside: where exactly does this paper sit relative to the more formal Genesis and Consciousness materials? The two bridge papers answer local parts of that question by clarifying scale-stable imprinting and stable-hull integration. The present paper answers the global version. It shows the whole chain at once.

The paper's governing sentence is therefore the following:

\begin{quote}
This paper does not prove the chain; it displays it, grades it, and marks the places where theorem-level results, programme-level bridges, interface claims, and structural resonances meet.
\end{quote}

That sentence expresses both the ambition and the restraint of the paper. It is ambitious because it attempts to present one coherent architecture from source-side emergence to finite self-description. It is restrained because it refuses to blur the difference between a theorem and a philosophical bridge.

The grading discipline is not arbitrary. It resonates with structural-realist programmes in philosophy of science, which separate structural content from extra-structural commitments \cite{Ladyman2007,French2014}. It also sits naturally beside Batterman's emphasis on asymptotic and cross-scale explanation, where explanatory power often outruns one-to-one microstructural reconstruction \cite{Batterman2002}. This paper extends that orientation by adding an explicit category of interface claims between mature vocabularies.

\section{The Full Imprint Chain}

The chain can be written in its shortest usable form:

\begin{quote}
source algebra $\to$ coarse-graining $\to$ scale-stable imprints $\to$ imprint-field $\to$ accessible stable hull $\to$ self-model $\to$ self-reading threshold $\to$ finite self-description horizon.
\end{quote}

This sequence is not meant as a single deduction. It is a structured path across multiple already developed works. Each arrow asks a different question:

\begin{center}
\begin{tabular}{P{0.34\linewidth}P{0.56\linewidth}}
\toprule
\textbf{Segment} & \textbf{Question} \\
\midrule
source algebra $\to$ coarse-graining & How does lower-level structure become effectively describable at all? \\
coarse-graining $\to$ scale-stable imprints & Which marks survive scale change as readable residues? \\
scale-stable imprints $\to$ imprint-field & How should emergence be read once its surviving residues are treated as trace-carrying? \\
imprint-field $\to$ accessible stable hull $\to$ self-model & How does a finite observer inherit only a structured subset of those traces as a world that can support self-location? \\
self-model $\to$ self-reading threshold & When does the model become reflexively available as itself imprint-bearing? \\
self-reading threshold $\to$ finite self-description horizon & Why must finite self-reading remain structurally incomplete? \\
\bottomrule
\end{tabular}
\end{center}

The paper's first task is simply to insist that these are not all the same kind of question. Confusion begins whenever they are treated as if they belonged to one epistemic register.

\section{Segment I: Algebraic Emergence}

The first segment concerns the move from source-side structure to coarse-grained observability. This is largely the territory of the Genesis Trilogy and of Foundation~I \cite{HAFFA,HAFFD,FoundationI}. The relevant lesson is not that every source-side detail survives. It is that emergence is already filtered by accessibility, projection, or effective selection.

This segment should be read as a mixture of formal and conditional work. In the stronger Genesis papers, coarse-grained observable structure is treated in explicit algebraic terms. In Foundation~I, the passage from source datum to accessibility remains conditional in precise ways. The point for the present map is that the chain does not begin with already-given minds or worlds. It begins with the problem of how any effective descriptive regime becomes available at all.

\section{Segment II: Scale-Stable Imprint Filtering}

The second segment is the task of the first bridge paper, \emph{Renormalization as Imprint Filtering}. Its role is to explain what stability means in the imprint vocabulary once scale change is taken seriously. The claim here is programme-level rather than theorem-level: renormalization language can be read as answering which readable residues survive coarse-graining.

This segment is important because without it the imprint vocabulary risks remaining too intuitive. ``Stable'' would then sound like a decorative adjective rather than a physically constrained notion. The bridge paper sharpens it by relating stability to relevance, universality, fixed-point behavior, and symmetry-breaking residue.

\section{Segment III: Imprint-Field and Universal Carrying}

The third segment is the conceptual core of \emph{The World as Imprint} \cite{WorldImprint}. Once scale-stable residues are in view, the emergent world can be read as an imprint-field: not a hidden substrate, but a domain of structures insofar as they bear readable traces of their generative constraints.

This segment is explicitly programme-level and philosophical. It is not a theorem, and it does not pretend to be one. Its contribution is to give a disciplined vocabulary for a structural fact already distributed across the more formal papers: emergence is not blank outcome but marked organization.

\section{Segment IV: Stable Hull Integration}

The fourth segment is the task of the second bridge paper, \emph{Stable Hulls as Imprint Integration}. Here the question is no longer merely which traces survive scale change, but which traces become a usable world for a finite observer. Foundation~II and the Brain Anchor paper already provide much of the architecture needed for this step \cite{FoundationII,BrainAnchor}. The bridge contribution is to recast that architecture in imprint language.

This segment is best understood as an interface claim. It connects a technical lineage about accessibility and stable hulls to a philosophical lineage about trace-carrying and self-reading. Its purpose is to make explicit the middle architecture between universal carrying and reflexive recognition.

The self-model literature is developed more directly in \emph{Stable Hulls as Imprint Integration} \S6; this roadmap only locates that engagement.

\section{Segment V: The Self-Reading Threshold}

The fifth segment returns to \emph{The World as Imprint}. There the central distinction is between an imprint-carrier and a self-reading system \cite{WorldImprint}. The self-reading threshold is crossed when a system can reflexively read its own world-model, or itself, as imprint-bearing.

This segment should again be read as programme-level rather than theorem-level. It provides a marker matrix and a conceptual threshold, not a full constitutive theory of first-person consciousness. Its value lies in holding apart several things often confused: structure, responsiveness, self-modeling, and reflexive self-reading.

\section{Segment VI: The Finite Horizon of Self-Description}

The final segment is developed most sharply in \emph{The Structural Horizon of Self-Description} \cite{PaperIII}. Once a system reaches organized self-reading, a further question appears: can it make itself fully transparent to itself? The answer given there is no. Finite self-description carries a structural horizon.

This segment occupies the strongest formal region of the later chain. It states not merely a philosophical caution but a principled bound on local self-description. In roadmap terms, it functions as the final restraint on any fantasy of complete reflexive closure.

\section{Epistemic Grading Map}

The chain becomes clearer once its internal grades are stated explicitly.

\begin{center}
\begin{tabular}{P{0.24\linewidth}P{0.33\linewidth}P{0.33\linewidth}}
\toprule
\textbf{Grade} & \textbf{Meaning} & \textbf{Examples in the chain} \\
\midrule
theorem-level / formal result & proved or formally developed within a technical framework with explicit conditions and inferential closure & local self-description bound, conditional hull construction, geometry/modular inseparability \\
programme-level bridge & structurally motivated and explanatorily useful, but not offered as an independent theorem in the target field & imprint principle, RG as imprint filtering, self-reading threshold as architectural criterion \\
interface claim & a claim that one mature language can be used to organize or reread another without reconstructing its full machinery & stable hulls as imprint integration, brain anchor as internal stable-hull mechanism \\
structural resonance / metaphor & reader-facing orientation not used as evidential premise or formal support & Buddhist resonance, Infinite Light metaphor, illustrative world-reading images \\
\bottomrule
\end{tabular}
\end{center}

This table is not optional. Without it, readers are likely to over-read weaker claims as if they were formal completions, or under-read stronger claims as if they were merely poetic. The map protects both sides: formal work is not diluted, and bridge work is not inflated. In particular, a programme-level bridge earns its place not by proving a new theorem, but by making a previously implicit research architecture explicit enough to guide later formalization and later reading.
Its core contribution is this grading map.

\section{How the Papers Fit Together}

The chain can now be mapped directly onto existing and planned papers.

\begin{center}
\begin{tabular}{P{0.34\linewidth}P{0.56\linewidth}}
\toprule
\textbf{Chain segment} & \textbf{Primary paper location} \\
\midrule
source algebra $\to$ coarse-graining & Genesis Trilogy; Foundation~I \\
coarse-graining $\to$ scale-stable imprints & \emph{Renormalization as Imprint Filtering} \\
scale-stable imprints $\to$ imprint-field & \emph{The World as Imprint} \\
imprint-field $\to$ accessible stable hull $\to$ self-model & \emph{Stable Hulls as Imprint Integration} \S4--\S6; Foundation~II; Brain Anchor \\
self-model $\to$ self-reading threshold & \emph{The World as Imprint} \\
self-reading threshold $\to$ finite self-description horizon & \emph{The Structural Horizon of Self-Description} \\
\bottomrule
\end{tabular}
\end{center}

The key point is that the bridge papers do not create a fourth series. They close local gaps between already existing clusters of work. Their natural place is interstitial rather than foundational in their own right. Papers 1 and 2 are therefore best treated as Zenodo working papers serving as interpretive companions to the present roadmap; they are not offered as independent technical contributions to renormalization theory or to cognitive science.

\section{Reading Paths for External Readers}

The body of work is now large enough that not every reader should be asked to read it all in chronological order. Three paths are therefore useful.

\subsection{Two-paper path}

Readers mainly interested in the philosophical architecture can begin with:

\begin{enumerate}
\item \emph{The World as Imprint}
\item \emph{The Imprint Chain}
\end{enumerate}

This path gives the overall frame and the internal grading without requiring the full technical background.

\subsection{Five-paper path}

Readers interested in structure and mind, but not in all of Genesis, can follow:

\begin{enumerate}
\item \emph{The World as Imprint}
\item \emph{Renormalization as Imprint Filtering}
\item \emph{Stable Hulls as Imprint Integration}
\item \emph{The Brain as a Plastic Observable-Algebra Network}
\item \emph{The Structural Horizon of Self-Description}
\end{enumerate}

This path captures the later half of the chain in a comparatively self-contained way.

\subsection{Full path}

Readers pursuing the complete architecture should read:

\begin{enumerate}
\item the Genesis Trilogy
\item the Consciousness Series
\item \emph{The World as Imprint}
\item the two bridge papers
\item the present roadmap
\end{enumerate}

The present paper is not the first paper for such readers. It is the paper that helps them not get lost once the archive becomes large.

\section{What the Chain Does Not Claim}

The map is useful only if it also states its non-claims clearly.

First, the chain is not all theorem-level. It contains formal results, bridges, interfaces, and resonances.

Second, the chain does not prove that physics validates Buddhism or any other contemplative tradition. Such resonances are explicitly reader-facing and non-evidential.

Third, the chain does not solve the hard problem. Even if self-reading architecture is clarified, why experience feels like something remains open.

Fourth, the chain does not establish soul survival, identity transfer, or complete self-transparency.

Fifth, the existence of a map does not mean the map is complete. The point is not closure by proclamation, but intelligibility by segmentation.

\section{Conclusion: Displayed, Not Proved}

The point of this paper has been organizational from the start. A large body of work can fail not because its parts are individually weak, but because readers cannot tell how the parts are supposed to hang together. The imprint chain is an attempt to prevent that failure.

What emerges from the map is a cleaner picture of the whole arc. The Genesis papers explain how effective structure emerges under accessibility and coarse-graining. The first bridge paper explains why some residues survive as scale-stable imprints. \emph{The World as Imprint} says how to read such a world philosophically. The second bridge paper explains how finite observers inherit only an integrated subset of that world as a usable hull. The later consciousness papers explain how such a regime can support self-modeling and why finite self-reading remains bounded.

This does not amount to a single theorem. It amounts to a graded route from emergence to self-reading. That is enough. A roadmap should not pretend to be the terrain. Its task is to make the terrain traversable.

\begin{thebibliography}{99}

\bibitem{WorldImprint}
S.~Liu,
\emph{The World as Imprint: From Trace-Carrying Structure to Self-Reading Mind},
Zenodo (2026), DOI: 10.5281/zenodo.20122718.

\bibitem{HAFFA}
S.~Liu,
\emph{Emergent Geometry from Coarse-Grained Observable Algebras: The Holographic Alaya-Field Framework},
Zenodo (2026), DOI: 10.5281/zenodo.18361706.

\bibitem{HAFFD}
S.~Liu,
\emph{Gravitational Phenomena as Emergent Properties of Observable Algebra Selection: A Structural Analysis},
Zenodo (2026), DOI: 10.5281/zenodo.18388881.

\bibitem{FoundationI}
S.~Liu,
\emph{Foundation I: From Source to Manifestation: A Conditional Same-Root Program for Geometry and Modular Flow},
Zenodo (2026), DOI: 10.5281/zenodo.19627786.

\bibitem{FoundationII}
S.~Liu,
\emph{Foundation II: Shared Manifestation: A Conditional Two-Source Program for Shared Accessible Sectors},
Zenodo (2026), DOI: 10.5281/zenodo.19641372.

\bibitem{FoundationIII}
S.~Liu,
\emph{Foundation III: The Perfuming Loop: A Conditional One-Step Source-Level Feedback Program},
Zenodo (2026), DOI: 10.5281/zenodo.19642189.

\bibitem{BrainAnchor}
S.~Liu,
\emph{The Brain as a Plastic Observable-Algebra Network: Internal Shared Stable Hulls, Neural Plasticity, and First-Person Report},
Zenodo (2026), DOI: 10.5281/zenodo.20044329.

\bibitem{PaperIII}
S.~Liu,
\emph{The Structural Horizon of Self-Description: Fisher Bounds, Small-Ball Remainders, and the Quantitative Explanatory Gap},
Consciousness Trilogy, Paper~III (2026), Zenodo, DOI: 10.5281/zenodo.19607724.

\bibitem{Ladyman2007}
J.~Ladyman and D.~Ross,
\emph{Every Thing Must Go: Metaphysics Naturalized},
Oxford University Press (2007).

\bibitem{French2014}
S.~French,
\emph{The Structure of the World: Metaphysics and Representation},
Oxford University Press (2014).

\bibitem{Batterman2002}
R.~W.~Batterman,
\emph{The Devil in the Details: Asymptotic Reasoning in Explanation, Reduction, and Emergence},
Oxford University Press (2002).

\end{thebibliography}

\end{document}
